\section{Strong chordality implies MAT-freeness}
\subsection{MAT-simplicial vertices and MAT-perfect elimination orderings}
\label{subsec:MAT-S-PEO}
 
The proof of the ``if'' part of our main Theorem \ref{thm:MAT-free-strong-chordal} (strong chordality implies MAT-freeness) requires more effort. 
We need a deeper understanding of the structure of graphs having MAT-labelings. 
In this subsection, we develop a fundamental study on such graphs analogous to the theory of (strongly) chordal graphs by introducing MAT- versions of simplicial vertex and perfect elimination ordering.
\begin{definition}[MAT-simplicial vertices]
\label{definition MAT-simplicial}
Given an edge-labeled graph $ (G,\lambda) $, a vertex $ v \in V_{G} $ is said to be \textbf{MAT-simplicial} if the following conditions hold. 
\begin{enumerate}[(MS1)]
\item\label{definition MAT-simplicial 1} $ v $ is a simplicial vertex of $ G $, that is, its neighborhood $ N_{G}(v) $ is a clique of $ G $.
\item\label{definition MAT-simplicial 2} $ \Set{\lambda(\{u,v\}) \in \mathbb{Z}_{>0} | u \in N_{G}(v) } = \{1,2, \dots, \deg_{G}(v)\} $, where $ \deg_{G}(v) =|N_{G}(v) |$ denotes the degree  of $ v $ in $ G $. 
\item\label{definition MAT-simplicial 3} For any distinct  $ u_{1}, u_{2} \in N_{G}(v) $, $ \lambda(\{u_{1}, u_{2}\}) < \max\{\lambda(\{u_{1}, v\}), \lambda(\{u_{2}, v\})\} $. 
\end{enumerate}
\end{definition}


Next we show the existence of MAT-simplicial vertices in the graphs having MAT-labelings.


\begin{lemma}
\label{lem:MATS-existence}
Let $ (G, \lambda) $ be an edge-labeled graph such that $ |V_{G}| \geq 2 $ and $ \lambda $ is an MAT-labeling of $ G $. 
\begin{enumerate}[(1)]
\item\label{MATS-existence 1} If $ G = K_{\ell} $ is a complete graph, then the endvertices of the edge with maximal label are MAT-simplicial. 
\item\label{MATS-existence 2} 
If $ G $ is noncomplete, then $ (G, \lambda) $ has two nonadjacent MAT-simplicial vertices. 
\end{enumerate}
\end{lemma}
\begin{proof}
First we prove part \ref{MATS-existence 1}. 
Let $ e_{0} = \{u_{0}, v_{0}\} \in E_{G} $ be the edge with maximal label. 
It suffices to prove that $ v_{0} $ is MAT-simplicial. 
First, \ref{definition MAT-simplicial 1} is clear. 
Next we show \ref{definition MAT-simplicial 2}. 
Note that $ \lambda|_{E_{G\setminus e_0}} $ is an MAT-labeling.
Then by Lemma \ref{restriction to intersection of maximal cliques}, the labelings~$ \lambda|_{E_{G[C]}} $ and $ \lambda|_{E_{G[X]}} $ are MAT-labelings, where $ C \coloneqq V_{G}\setminus\{u_{0}\} $ and $ X \coloneqq V_{G} \setminus \{u_{0}, v_{0}\} $. 
Comparing the exponents of $ \mathcal{A}_{G[C]} $ and $ \mathcal{A}_{G[X]} $, we have $ \Set{\lambda(\{v, v_{0}\}) \in \mathbb{Z}_{>0} | v \in X } = [\ell -2] $. 
Since $ \lambda(e_{0}) = \ell -1 $, $ \Set{\lambda(\{v, v_{0}\}) \in \mathbb{Z}_{>0} | v \in N_{G}(v_{0}) } = [\ell -1] $.
Thus \ref{definition MAT-simplicial 2} is satisfied. 

Lastly,  we show \ref{definition MAT-simplicial 3}. 
Let $ u,v \in X $ and write $ \lambda(\{u,v\}) = k $. 
We want to show~$ k<\max\{\lambda(\{u,v_{0}\}), \lambda(\{v,v_{0}\})\} $. 
Since both $ \lambda|_{E_{G[C]}} $ and $ \lambda|_{E_{G[X]}} $ are MAT-labelings, \ref{definition MAT-labeling triangle} implies $ k \leq \max\{\lambda(\{u,v_{0}\}), \lambda(\{v,v_{0}\})\} $. 
If the equality happens, then it contradicts to \ref{definition MAT-labeling closure} of $ \lambda|_{E_{G[C]}} $. 
Therefore $ k < \max\{\lambda(\{u,v_{0}\}), \lambda(\{v,v_{0}\})\} $. 
Moreover, $ \max\Set{\lambda(\{v,u_{0}\}) | v \in X} = \ell -2 < \ell - 1 = \lambda(e_{0}) $, since $ \lambda|_{E_{G\setminus v_{0}}} $ is also an MAT-labeling by Lemma \ref{restriction to intersection of maximal cliques}. 
Therefore \ref{definition MAT-simplicial 3} holds. 
Thus $ v_0 $ is MAT-simplicial. 

Now we prove part \ref{MATS-existence 2}. 
 We proceed induction on $ \ell = |V_{G}| $. 
If $ \ell=2 $, then the assertion holds trivially. 
Suppose $ \ell \geq 3 $. 
We may assume that $ G $ is connected. 
So let~$ a,b \in V_{G} $ be nonadjacent vertices and $ S $ a minimal $ (a,b) $-separator. 
Then $ S $ is a clique by Theorem \ref{Dirac minimal vertex separator} and $ S \in \mathcal{P}_{G} $ by Remark \ref{Ho and Lee}. 
Hence $ \lambda|_{E_{G[S]}} $ is an MAT-labeling of~$ G[S] $ by Lemma \ref{restriction to intersection of maximal cliques}.

Let $ A $ be the vertex set of the connected component containing $ a $ of $ G\setminus S $ and $ B \coloneqq V_{G} \setminus (A \cup S) $. 
We will show that $ \lambda|_{E_{G[A\cup S]}} $ and $ \lambda|_{E_{G[B\cup S]}} $ are MAT-labelings. 
By Proposition \ref{restriction to subset}, it suffices to show \ref{definition MAT-labeling triangle}. 
Let $ e \in \pi_{k} \cap E_{G[A \cup S]} $. 
Then $ e $ forms at most $ k-1 $ triangles with edges in $ E_{k-1} \cap E_{G[A \cup S]} $ since $ \lambda $ is an MAT-labeling. 
When~$ e \in E_{G[S]} $, $ e $ forms exactly $ k-1 $ triangles with edges in $ E_{k-1} \cap E_{G[S]} $. 
Therefore $ e $ forms exactly $ k-1 $ triangles with edges in $ E_{k-1} \cap E_{G[A \cup S]} $. 
Suppose that at least one endvertex of $ e $ belongs to $ A $. 
Then $ e $ cannot form a triangle with a vertex in $ B $. 
Hence $ e $ forms exactly $ k-1 $ triangles with edges in $ E_{k-1} \cap E_{G[A \cup S]} $. 
Thus $ \lambda|_{E_{G[A\cup S]}} $ is an MAT-labeling.
By a similar way, one can prove that $ \lambda|_{E_{G[B\cup S]}} $ is an MAT-labeling.


Next we show that $ A \setminus S $ contains an MAT-simplicial vertex of $ G $. 
Note that if~$ v \in A \setminus S $ is MAT-simplicial in $ G[A\cup S] $, then $ v $ is also MAT-simplicial in $ G $ since $ N_{G}(v) \subseteq A \cup S $. 
If $ G[A \cup S] $ is a complete graph, then the endvertices of the edge $ e_{0} $ in~$ G[A \cup S] $ with maximal label is MAT-simplicial in~$ G[A \cup S] $ by  part \ref{MATS-existence 1}. 
Since $ \lambda|_{E_{G[S]}} $ is an MAT-labeling, at least one endvertex of $ e_{0} $ belongs to $ A \setminus S $ by Proposition \ref{MAT-labeling on complete graph}, which is a desired MAT-simplicial vertex. 
If $ G[A \cup S] $ is not a complete graph, then by the induction hypothesis $ G[A \cup S] $ has two nonadjacent MAT-simplicial vertices. 
At least one of them belongs to $ A \setminus S $ since $ S $ is a clique. 
Thus $ A \setminus S $ contains an MAT-simplicial vertex of $ G $. 

Similarly, $ B \setminus S $ contains an MAT-simplicial vertex of $ G $. 
Therefore $ G $ has two nonadjacent MAT-simplicial vertices. 
\end{proof}




The following is a first important property of MAT-simplicial vertices.
\begin{proposition}\label{MAT-simplicial}
Let $ (G, \lambda) $ be an edge-labeled graph with $ |V_{G}| \geq 2 $. 
Suppose that $ v \in V_{G} $ is an MAT-simplicial vertex of $ (G,\lambda) $. 
The following are equivalent. 
\begin{enumerate}[(1)]
\item $ \lambda $ is an MAT-labeling of $ G $. 
\item $ \lambda|_{E_{G\setminus v}} $ is an MAT-labeling of $ G\setminus v $. 
\end{enumerate}
\end{proposition}
\begin{proof}
First we prove $ (1) \Rightarrow (2) $. 
Let $ \pi_{k}^{\prime} \coloneqq (\lambda|_{E_{G\setminus v}})^{-1}(k) =\pi_k \cap E_{G\setminus v}$ and $ E_{k-1}^{\prime} \coloneqq \pi^{\prime}_{1} \sqcup \dots \sqcup \pi^{\prime}_{k-1} $ for $ k \in \mathbb{Z}_{>0} $. 
By Proposition  \ref{restriction to subset}, we only need to prove \ref{definition MAT-labeling triangle} of~$ \lambda|_{E_{G\setminus v}} $.

Let $ e \in \pi_{k}^{\prime} $. 
Since $ e \in \pi_{k}^{\prime} \subseteq \pi_{k} $, $ e $ forms exactly $ k-1 $ triangles with edges in $ E_{k-1} $. 
These triangles do not contain the vertex $ v $ because by \ref{definition MAT-simplicial 3} the number of edges incident to $ v $ with label less than $ k $ is at most $ 1 $. 
Therefore $ e $ forms exactly $ k-1 $ triangles with edges in $ E_{k-1}^{\prime} $. 
Thus $ \lambda|_{E_{G\setminus v}} $ is an MAT-labeling of $ G\setminus v $. 

Next we prove $ (2) \Rightarrow (1) $. 
Let $ k \in \mathbb{Z}_{>0} $.
By  \ref{definition MAT-simplicial 2}, $ v $ is a leaf or  an isolated vertex of $ \pi_{k} $. 
Moreover, since $ \pi_{k}^{\prime} $ is a forest, $ \pi_{k} $ is a forest. 
This shows \ref{definition MAT-labeling forest}. 

To show \ref{definition MAT-labeling closure}, suppose $ \cl_{G}(\pi_{k}) \cap E_{k-1} \neq \varnothing $ and take $ e \in \cl_{G}(\pi_{k}) \cap E_{k-1} $. 
Then there exists a cycle $ C $ in $G$ such that $ e \in C $ and $ C \setminus e \subseteq \pi_{k} $. 
If $ e $ is not incident to~$ v $ (in particular, $v$ is not a vertex of $C$), then $ e \in E_{k-1}^{\prime} $ and $ C \setminus e \subseteq \pi_{k}^{\prime} $. 
Therefore $ e \in \cl_{G\setminus v}(\pi_{k}^{\prime}) \cap E_{k-1}^{\prime} = \varnothing$, a contradiction. 
Hence $ e $ is incident to $ v $, and $ C $ contains an edge $ \{v,w\} $ with $ \lambda(\{v,w\}) = k $. 
Write $ e = \{u,v\} $. 
Then $ \{u,w\} \in E_{G} $ by \ref{definition MAT-simplicial 1} and $ \lambda(\{u,w\}) < \max\{\lambda(\{u,v\}), \lambda(\{v,w\})\} = k $ by \ref{definition MAT-simplicial 3}. 
Hence $ \{u,w\} \in E_{k-1}^{\prime} $ (in particular, $C$ has length at least $4$).
Moreover,  $ \{u,w\}$ and $ C \setminus\{\{v,w\}, e\} \subseteq \pi_{k}^{\prime} $ form a cycle and hence $ \{u,w\} \in \cl_{G\setminus v}(\pi_{k}^{\prime}) \cap E_{k-1}^{\prime}= \varnothing$, a contradiction. 

Finally, we prove \ref{definition MAT-labeling triangle}. 
Let $ e \in \pi_{k} $.
If $ e \in \pi_{k}^{\prime} $ (i.e, $ e $ is not incident to $ v $), then~$ e $ forms exactly $ k-1 $ triangles with some edges in $ E_{k-1}^{\prime} $. 
If one endvertex of $ e $ is not adjacent to $ v $, then $ e $ and $v$ cannot form a triangle. 
If both endvertices of $ e $ are adjacent to $ v $, then at least one edge of the triangle containing $ e $ and $ v $ has label greater than $ k $ by \ref{definition MAT-simplicial 3}. 
In either case, $ e $ forms exactly $ k-1 $ triangles with some edges in $ E_{k-1} $. 
Now consider  $ e \in \pi_{k} \setminus \pi_{k}^{\prime} $ (\ie  $ e $ is incident to $ v $). 
By  \ref{definition MAT-simplicial 2}, we can index the elements of~$ N_{G}(v)$  as $\{u_{1}, \dots, u_{d}\} $  where $k\le d=\deg_G(v)$ so that $ e = \{u_{k}, v\} $ and $ \lambda(\{u_{i}, v\}) = i $ for every $ i \in [d] $. 
Hence $ e $ forms exactly $ k-1 $ triangles given by $ \{u_{i}, u_{k}, v\} $ for $ i \in [k-1] $ with some edges in $ E_{k-1} $.  
Thus $ \lambda $ is an MAT-labeling.  
\end{proof}

\begin{definition}[MAT-PEO]
Given an edge-labeled graph $ (G,\lambda) $, an ordering $ (v_{1}, \dots, v_{\ell}) $ of $ G $ is said to be a \textbf{MAT-perfect elimination ordering (MAT-PEO)} of $ (G,\lambda) $ if $ v_{i} $ is MAT-simplicial in $ (G_{i}, \lambda_{i}) $ for each $ i \in [\ell] $, where $ G_{i} \coloneqq G[\{v_{1}, \dots, v_{i}\}] $ and $ \lambda_{i} \coloneqq \lambda|_{E_{G_{i}}} $. 
\end{definition}

In particular, any MAT-PEO is a PEO.
The theorem below exhibits a strong connection between MAT-labelings and MAT-PEOs which can be seen as an analogue of Theorem \ref{thm:chordal-PEO}.
\begin{theorem}[MAT-labelings and MAT-PEOs]
\label{MAT-labeling ordering}
Given an edge-labeled graph $ (G, \lambda) $, the following are equivalent. 
\begin{enumerate}[(1)]
\item $ \lambda $ is an MAT-labeling of $ G $. 
\item There exists an MAT-PEO of $ (G, \lambda) $. 
\end{enumerate}
\end{theorem}
\begin{proof}
Both implications can be proved by induction on $ \ell = |V_{G}| $.
The implication $ (2) \Rightarrow (1) $ is easy thanks to Proposition \ref{MAT-simplicial}. 
Let us show the converse (which is also not hard). 
When $ \ell = 1 $, the assertion is true.
Suppose $ \ell \geq 2 $. 
By Lemma \ref{lem:MATS-existence}, there exists an MAT-simplicial vertex $ v_{\ell} $ of $ (G, \lambda) $. 
Then $ \lambda|_{E_{G\setminus v_{\ell}}} $ is an MAT-labeling of~$ G\setminus v_{\ell} $ by Proposition \ref{MAT-simplicial}. 
By the induction hypothesis, $ (G\setminus v_{\ell}, \lambda|_{E_{G\setminus v_{\ell}}}) $ has an MAT-PEO $ (v_{1}, \dots, v_{\ell -1}) $. 
Thus $ (v_{1}, \dots, v_{\ell}) $ is an MAT-PEO of $ (G,\lambda) $. 
\end{proof}

We complete this subsection by giving two lemmas on (extensions of) MAT-labelings and MAT-PEOs of complete graphs. 
MAT-labelings of complete graphs will play a crucial role in the next subsection. 
